% Generated draft: requires mathematical and source review.
\documentclass[11pt]{article}
\usepackage[margin=25mm]{geometry}
\usepackage{fontspec}
\setmainfont{Latin Modern Roman}
\usepackage{amsmath,amssymb,mathtools,unicode-math}
\setmathfont{Latin Modern Math}
\usepackage{enumitem,needspace,xurl,hyperref}
\setlength{\parindent}{0pt}
\setlength{\parskip}{5pt}
\setlength{\emergencystretch}{3em}
\newcommand{\N}{\mathbb{N}}
\newcommand{\Z}{\mathbb{Z}}
\newcommand{\Q}{\mathbb{Q}}
\newcommand{\R}{\mathbb{R}}
\newcommand{\C}{\mathbb{C}}
\newcommand{\F}{\mathbb{F}}
\newcommand{\A}{\mathbb{A}}
\newcommand{\PP}{\mathbb{P}}
\newcommand{\E}{\mathbb{E}}
\newcommand{\eps}{\varepsilon}
\newcommand{\dd}{\,\mathrm{d}}
\begin{document}
\section*{UnsolvedMath ID: 30006227}
\textbf{Source problem code:} catalogue entry
\subsection{Definitions and mathematical statement}
\textbf{Unrecoverable statement.} The supplied record for permanent UnsolvedMath ID \(30006227\), source code \texttt{catalogue entry}, contains no recoverable mathematical assertion. Its extracted text begins with ``Conjecture'' followed by a dotted leader and page number \(586\), then continues with apparent table-of-contents entries for other contributions. No mathematical objects, definitions, hypotheses, domains, boundary conditions, quantifiers, or conclusion of the intended conjecture are supplied. Consequently, a precise mathematical statement cannot be formulated from this evidence.

The registry title, ``Landis-Conjecture Table-of-Contents Fragment'', does not specify a version of the Landis conjecture or establish that any particular formulation is intended. The inequality \(1<p<2\) occurs only in an adjacent contribution title; it cannot be assigned as a hypothesis of the missing conjecture. Likewise, the adjacent contributors cannot be identified as its proposers.

The supplied statement status is \texttt{unrecoverable}. The supplied verification note reports that the official report confirms the extraction is table-of-contents debris rather than a conjecture statement; this verification claim has not been independently checked here. The supplied \texttt{open} label is a curation classification adopted because the truncated extraction cannot be matched to a proof or counterexample, and does not establish that an identifiable mathematical conjecture remains open.
\subsection{Short English statement}
The supplied text is a table-of-contents fragment, so no precise conjecture can be recovered for UnsolvedMath ID \(30006227\). Its title mentions the Landis conjecture, but supplies no formulation. The recorded open status reflects supplied curation of an unidentified assertion. Recovery of the intended passage and human mathematical review are required.
\subsection{Sources}
\begin{itemize}
\item Local UnsolvedMath registry, supplied complete record: permanent ID \(30006227\); exact source code \texttt{catalogue entry}; collection ``Oberwolfach Reports Open Problems''; proposed year \(2025\). No proposer is identified in the record.
\item Supplied report citation: \emph{Mini-Workshop: Hardy Inequalities in Discrete and Continuum Settings} (2025), Oberwolfach Reports, pp.~579--614, DOI \texttt{10.4171/owr/2025/14}. Supplied source URL and literature-source link, labelled ``Oberwolfach Report 14/2025'': \url{https://doi.org/10.4171/owr/2025/14}. Oberwolfach Reports are described in the supplied background as documenting mathematical workshops and their open-problem discussions. This report citation establishes supplied provenance, but the record provides no usable passage specifying the intended conjecture.
\item Supplied dated literature curation: \texttt{literature\_checked\_at} is exactly \texttt{2026-08-22}. Its assessment says that the extraction is truncated to the word ``Conjecture'', contains no identifiable assertion that can be matched to a proof or counterexample, and is therefore classified as open under the curation's required normalization. The report above is the sole listed checked source. This is supplied dated curation, not an independently verified literature assessment; the supplied background explicitly calls for independent verification of time-sensitive status and cited claims.
\item Supplied statement-verification note: the official report is said to confirm that the extraction is table-of-contents debris, not a conjecture statement. The record marks the statement \texttt{unrecoverable}; no independent source visit or verification is asserted here.
\item Context retained in the supplied extraction: ``Conjecture'', p.~586; Baptiste Devyver (joint with Louis Dupaigne and Pierre-Damien Thizy), ``Sharp Hardy-Poincaré-Sobolev inequalities in the Caffarelli-Kohn-Nirenberg hyperbolic space'', p.~587; Bartlomiej Dyda (joint with Sven Jarohs, Firoj Sk; Michal Kijaczko), ``Logarithmic Hardy inequality. Fractional Sobolev--Hardy--Maz'ya inequality for \(1<p<2\)'', p.~590. These entries provide context only and do not supply the missing assertion or its attribution.
\item Further contextual entries: Luca Fanelli, ``Uncertainty Principles: from Hardy to Resolvent Inequalities'', p.~592; Florian Fischer (joint with Christian Rose), ``Optimal Poincaré-Hardy Inequalities on Manifolds and Graphs'', p.~593; Rupert L. Frank (joint with Ari Laptev and Timo Weidl), ``Some remarks on one-dimensional Hardy inequalities'', p.~595; Gabriele Grillo (joint with Elvise Berchio, Dario D. Monticelli, Matteo Muratori and Fabio Punzo), ``Fractional nonlinear diffusions on manifolds: well-posedness and smoothing effects'', p.~597; Shubham Gupta, ``Hardy inequalities on integer lattices'', p.~598.
\item Remaining contextual entries: Hynek Kovařík (joint with Luca Fanelli), ``Quantitative Hardy inequalities for magnetic Hamiltonians'', p.~600; David Krejčiřík (joint with Borbala Gerhat and František Štampach), ``Criticality transition for powers of the discrete Laplacian'', p.~601; Felix Pogorzelski (joint with Philipp Hake, Matthias Keller and Yehuda Pinchover), ``Discrete optimal Hardy type inequalities and lattices in H-type groups'', p.~603; Prasun Roychowdhury (joint with Debdip Ganguly and Jotsaroop Kaur), ``Hardy and Rellich identities for Grushin operators'', p.~606. None is identified in the supplied record as the missing conjecture's author.
\end{itemize}
\end{document}
